\documentclass[11pt]{article}

% Appendix to Lecture 4: measure-theoretic probability
\usepackage[margin=1in]{geometry}
\usepackage[T1]{fontenc}
\usepackage{lmodern}
\usepackage{amsmath,amssymb,amsthm}
\usepackage{enumitem}
\usepackage{xr-hyper}
\usepackage{hyperref}
% Theorem numbers of Lectures 3 and 4 are read from lecture3.aux and lecture4.aux,
% so compile lecture3.tex and lecture4.tex first.
\externaldocument[L3-]{../Lecture3/lecture3}[../Lecture3/lecture3.pdf]
\externaldocument[L4-]{lecture4}[lecture4.pdf]
\usepackage{xcolor}
\usepackage{listings}

% Python code in the programming exercises.
\lstset{
    language=Python,
    basicstyle=\ttfamily\small,
    keywordstyle=\color{blue!70!black},
    commentstyle=\color{gray!70!black}\itshape,
    stringstyle=\color{orange!70!black},
    showstringspaces=false,
    upquote=true,
    columns=fullflexible,
    keepspaces=true,
    frame=single,
    rulecolor=\color{gray!40},
    backgroundcolor=\color{gray!6},
    xleftmargin=0.5em,
    xrightmargin=0.5em,
}
% Inline Python code, e.g. \py{np.linalg.solve(A, b)}; unlike \lstinline, it can break at spaces.
\newcommand{\py}[1]{\texttt{\detokenize{#1}}}

% Norms and inner products.
\newcommand{\norm}[1]{\lVert #1 \rVert}
\newcommand{\ip}[2]{\langle #1, #2 \rangle}

% Probability.
\newcommand{\E}{\mathbb{E}}
\DeclareMathOperator{\Var}{Var}
\DeclareMathOperator{\Cov}{Cov}
\DeclareMathOperator{\KL}{KL}

\title{Lecture 4, Appendix: Measure-Theoretic Probability}
\author{Math for Machine Learning}
\date{\today}

% The sections of the appendix are numbered B.1, B.2, ... (the appendix to Lecture 3 is A);
% widen the number boxes in the table of contents.
\renewcommand{\thesection}{B.\arabic{section}}
\usepackage{tocloft}
\setlength{\cftsecnumwidth}{2.9em}
\setlength{\cftsubsecnumwidth}{3.2em}

\theoremstyle{definition}
\newtheorem{definition}{Definition}[section]
\newtheorem{example}[definition]{Example}
\newtheorem{exercise}[definition]{Exercise}
\theoremstyle{plain}
\newtheorem{theorem}[definition]{Theorem}
\newtheorem{proposition}[definition]{Proposition}
\newtheorem{lemma}[definition]{Lemma}
\newtheorem{corollary}[definition]{Corollary}
\theoremstyle{remark}
\newtheorem{remark}[definition]{Remark}

\begin{document}

\maketitle
\tableofcontents

\section*{Introduction}
\addcontentsline{toc}{section}{Introduction}

In Lecture~4, we assigned probabilities to events and computed expectations as sums and integrals, without asking which sets of outcomes are events (Remark~\ref{L4-rem:measure-theory} of Lecture~4). This is enough for discrete random variables and for random variables with densities. But some distributions are neither discrete nor continuous, such as the output of a ReLU unit with a Gaussian input, and some questions are about infinitely many random variables at once, such as whether the sample means converge along a sequence of observations. Measure theory treats all of them in a single framework: events form a $\sigma$-algebra, probabilities are measures, random variables are measurable functions, and expectations are Lebesgue integrals, with powerful theorems for exchanging limits and expectations. This appendix builds probability theory on this foundation.
\begin{itemize}
    \item Sections~\ref{sec:measures}--\ref{sec:lebesgue} introduce $\sigma$-algebras and measures, random variables as measurable functions, and the Lebesgue integral.
    \item Sections~\ref{sec:convergence-theorems}--\ref{sec:product} prove the convergence theorems for integrals, and discuss densities with the Radon--Nikodym theorem, and product measures with the Fubini--Tonelli theorem.
    \item Sections~\ref{sec:conditional-expectation} and~\ref{sec:martingales} define the conditional expectation given a $\sigma$-algebra, as an orthogonal projection, and study martingales, the mathematical model of fair games.
    \item Section~\ref{sec:convergence} compares the different ways in which random variables converge, and proves the strong law of large numbers.
\end{itemize}
We use the notation of Lecture~4. Some proofs use Hilbert spaces, which are treated in the appendix to Lecture~3.

\section{Measures and Probability Spaces}\label{sec:measures}

A random experiment has a set $\Omega$ of possible outcomes, the sample space, and we want to assign probabilities to sets of outcomes, the events. If there are only finitely or countably many outcomes, every set of outcomes can be an event. For experiments with a continuum of outcomes, such as a uniformly random point of $[0, 1]$, this is impossible (Remark~\ref{rem:vitali}), and we have to specify which sets are events. Measure theory provides the framework: the events form a $\sigma$-algebra, and the probability is a measure on it.

\begin{definition}[$\sigma$-algebra]\label{def:sigma-algebra}
    A $\sigma$-algebra on a set $\Omega$ is a collection $\mathcal{F}$ of subsets of $\Omega$ such that
    \begin{enumerate}
        \item $\Omega \in \mathcal{F}$;
        \item if $A \in \mathcal{F}$, then $A^c = \Omega \setminus A \in \mathcal{F}$;
        \item if $A_1, A_2, \ldots \in \mathcal{F}$, then $A_1 \cup A_2 \cup \cdots \in \mathcal{F}$.
    \end{enumerate}
    The pair $(\Omega, \mathcal{F})$ is called a measurable space, and the sets in $\mathcal{F}$ are called measurable sets, or events.
\end{definition}

A $\sigma$-algebra contains $\emptyset = \Omega^c$, and since $\bigcap_k A_k = \big(\bigcup_k A_k^c\big)^c$ (De Morgan's law), it is also closed under countable intersections, and under finite unions, intersections and differences $A \setminus B = A \cap B^c$. So the usual set operations, applied countably many times, lead from events to events. The ``$\sigma$'' stands for countable unions; closedness under countable, not only finite, operations is what allows us to take limits.

\begin{example}[$\sigma$-algebras]\label{ex:sigma-algebras}
    \begin{enumerate}
        \item The power set $2^\Omega$, the set of all subsets of $\Omega$, is a $\sigma$-algebra; we use it whenever $\Omega$ is countable. The smallest $\sigma$-algebra on $\Omega$ is $\{\emptyset, \Omega\}$.
        \item (Information.) Let $\Omega$ be partitioned into disjoint sets $B_1, \ldots, B_n$. The unions of some of the $B_i$ form a $\sigma$-algebra $\mathcal{G}$ with $2^n$ elements. It describes the information of an observer who only learns which $B_i$ contains the outcome: an event belongs to $\mathcal{G}$ if and only if this observer can decide whether it has happened. For two dice, an observer who only learns the sum knows whether the event ``the sum is even'' has happened, but not whether ``the first die shows $6$'' has. We use $\sigma$-algebras in this way to describe information in Sections~\ref{sec:conditional-expectation} and~\ref{sec:martingales}.
        \item (Generated $\sigma$-algebras.) An intersection of $\sigma$-algebras on $\Omega$ is again a $\sigma$-algebra. So for every collection $\mathcal{C}$ of subsets of $\Omega$, there is a smallest $\sigma$-algebra $\sigma(\mathcal{C})$ that contains $\mathcal{C}$, namely the intersection of all $\sigma$-algebras that contain $\mathcal{C}$. It is called the $\sigma$-algebra generated by $\mathcal{C}$.
        \item (Borel sets.) The Borel $\sigma$-algebra $\mathcal{B}(\mathbb{R}^n)$ is the $\sigma$-algebra generated by the open subsets of $\mathbb{R}^n$, and its elements are called Borel sets. They include all open and closed sets, all boxes and all countable sets, and everything that can be built from them by countably many set operations. The Borel $\sigma$-algebra is also generated by the boxes $(a_1, b_1] \times \cdots \times (a_n, b_n]$, and for $n = 1$ by the intervals $(-\infty, x]$ with $x \in \mathbb{R}$.
    \end{enumerate}
\end{example}

\begin{definition}[Measure and probability space]\label{def:measure}
    A measure on a measurable space $(\Omega, \mathcal{F})$ is a function $\mu: \mathcal{F} \to [0, \infty]$ with $\mu(\emptyset) = 0$ that is countably additive: if $A_1, A_2, \ldots \in \mathcal{F}$ are pairwise disjoint, then
    \[\mu(A_1 \cup A_2 \cup \cdots) = \mu(A_1) + \mu(A_2) + \cdots.\]
    The triple $(\Omega, \mathcal{F}, \mu)$ is called a measure space. A probability measure is a measure $P$ with $P(\Omega) = 1$, and a probability space is a measure space $(\Omega, \mathcal{F}, P)$ with a probability measure $P$.
\end{definition}

A probability measure satisfies the axioms of Definition~\ref{L4-def:probability} of Lecture~4, except that only the sets in $\mathcal{F}$ are events. These are Kolmogorov's axioms of probability (1933). A measure is finite if $\mu(\Omega) < \infty$, and $\sigma$-finite if $\Omega$ is a countable union of sets of finite measure.

\begin{example}[Measures]\label{ex:measures}
    \begin{enumerate}
        \item (Discrete probability.) Let $\Omega = \{\omega_1, \omega_2, \ldots\}$ be countable, and let $p_1, p_2, \ldots \geq 0$ with $\sum_i p_i = 1$. Then $P(A) = \sum_{i : \omega_i \in A} p_i$ defines a probability measure on $2^\Omega$; it is countably additive, since series of nonnegative terms can be rearranged. If $\Omega$ is finite and all outcomes are equally likely, then $P(A) = \frac{|A|}{|\Omega|}$.
        \item (Counting and Dirac measures.) The counting measure $\mu(A) = |A|$, the number of elements of $A$ (or $\infty$), is a $\sigma$-finite measure on $2^{\mathbb{Z}}$. The Dirac measure $\delta_x(A) = \mathbf{1}_A(x)$ is the probability measure of an outcome that is $x$ with certainty.
        \item (Lebesgue measure.) There is exactly one measure $\lambda$ on $\mathcal{B}(\mathbb{R}^n)$ that assigns to every box its volume. We do not prove this theorem; it follows from Carath\'eodory's extension theorem, which extends a function on boxes to the $\sigma$-algebra that they generate. Since $\lambda(\{x\}) = 0$, countable additivity gives $\lambda(A) = 0$ for every countable set $A$, for example $\lambda(\mathbb{Q}) = 0$. The uniform distribution on $[0, 1]$ is the probability measure $P(A) = \lambda(A)$ on the Borel subsets of $[0, 1]$.
    \end{enumerate}
\end{example}

\begin{remark}[Why not all subsets?]\label{rem:vitali}
    There is no measure on all subsets of $[0, 1)$ that gives every interval its length and is invariant under translations modulo $1$. Here is the idea of Vitali's argument (1905). Call $x, y \in [0, 1)$ equivalent if $x - y$ is rational, and choose one point from each equivalence class; let $V$ be the set of chosen points (this needs the axiom of choice). The translates $V_q = \{v + q \bmod 1 : v \in V\}$ for the countably many rationals $q \in [0, 1)$ are pairwise disjoint and cover $[0, 1)$. If $V$ had a measure $m$, then so would every $V_q$, and countable additivity would give $1 = m + m + m + \cdots$, which is impossible both for $m = 0$ and for $m > 0$. Sets like $V$ never arise from explicit constructions; all sets that we meet in practice are Borel sets.
\end{remark}

The rules of Proposition~\ref{L4-prop:probability-rules} of Lecture~4 hold for the events of every probability space, with the same proofs, and two of them, monotonicity and the union bound, hold for every measure. Measures have one more property, which allows us to take limits.

\begin{proposition}[Continuity]\label{prop:continuity}
    Let $(\Omega, \mathcal{F}, \mu)$ be a measure space, and let $A_1, A_2, \ldots \in \mathcal{F}$.
    \begin{enumerate}
        \item If $A_1 \subseteq A_2 \subseteq \cdots$, then $\mu\big(\bigcup_k A_k\big) = \lim_{k \to \infty} \mu(A_k)$.
        \item If $A_1 \supseteq A_2 \supseteq \cdots$ and $\mu(A_1) < \infty$, then $\mu\big(\bigcap_k A_k\big) = \lim_{k \to \infty} \mu(A_k)$.
    \end{enumerate}
\end{proposition}

\begin{proof}
    (1) Let $A_0 = \emptyset$ and $B_k = A_k \setminus A_{k-1}$. The $B_k$ are disjoint, $A_n = B_1 \cup \cdots \cup B_n$, and the $B_k$ have the same union as the $A_k$, so $\mu\big(\bigcup_k A_k\big) = \sum_k \mu(B_k) = \lim_{n \to \infty}\sum_{k \leq n} \mu(B_k) = \lim_{n \to \infty} \mu(A_n)$. (2) The sets $A_1 \setminus A_k$ increase to $A_1 \setminus \bigcap_k A_k$, and $\mu(A_1 \setminus A_k) = \mu(A_1) - \mu(A_k)$, since $\mu(A_1) < \infty$. Now apply (1).
\end{proof}

Without the condition $\mu(A_1) < \infty$, part (2) can fail: the intervals $[k, \infty)$ decrease to $\emptyset$, but all have infinite Lebesgue measure.

An event with probability $1$ happens almost surely (a.s.), and an event with probability $0$ is a null event. Null events need not be impossible: a uniformly random point of $[0, 1]$ is irrational almost surely, since $\lambda(\mathbb{Q} \cap [0, 1]) = 0$, although every rational number in $[0, 1]$ is a possible outcome.

\begin{definition}[Independence]\label{def:independence}
    Infinitely many events are independent if any finitely many of them are, in the sense of Definition~\ref{L4-def:independence} of Lecture~4. More generally, $\sigma$-algebras $\mathcal{G}_1, \ldots, \mathcal{G}_n \subseteq \mathcal{F}$ are independent if $P(G_1 \cap \cdots \cap G_n) = P(G_1)\cdots P(G_n)$ for all $G_1 \in \mathcal{G}_1, \ldots, G_n \in \mathcal{G}_n$: they describe pieces of information that tell us nothing about each other.
\end{definition}

\begin{example}[Infinitely many coin flips]\label{ex:coin-flips}
    Many questions about random sequences need infinitely many independent events on one probability space. Let $\Omega = [0, 1)$ with the uniform distribution, and let $D_1(\omega), D_2(\omega), \ldots \in \{0, 1\}$ be the digits of the binary expansion $\omega = \sum_k D_k(\omega)2^{-k}$ (choosing the expansion that does not end with infinitely many $1$s). For every pattern $d_1, \ldots, d_n$, the set of $\omega$ whose first $n$ digits are $d_1, \ldots, d_n$ is an interval of length $2^{-n}$, so $P(D_1 = d_1, \ldots, D_n = d_n) = 2^{-n}$. Summing over the digits that we do not care about, it follows that the events $\{D_k = 1\}$ are independent, with probability $\frac{1}{2}$ each: the digits behave exactly like independent fair coin flips, with $1$ for heads. The events ``the first $n$ flips are tails'' have the probabilities $2^{-n}$ and decrease to the event ``all flips are tails'', so by continuity, its probability is $\lim_n 2^{-n} = 0$: almost surely, heads eventually appears. (The event ``all flips are tails'' is $\{0\}$, a possible but null event.) In the same way, one can construct a probability space that carries an infinite sequence of independent random variables (Definition~\ref{def:independent-rv}) with any given distributions; this is a theorem of Kolmogorov, which we do not prove.
\end{example}

\section{Random Variables and Their Distributions}\label{sec:random-variables}

A random variable is a numerical quantity that depends on the outcome of the experiment. To ask for the probability that it lies in a set $B$, the set of outcomes for which it does must be an event.

\begin{definition}[Random variable and distribution]\label{def:random-variable}
    Let $(\Omega, \mathcal{F}, P)$ be a probability space. A random variable is a measurable function $X: \Omega \to \mathbb{R}$, i.e., a function with
    \[\{X \in B\} = \{\omega \in \Omega : X(\omega) \in B\} \in \mathcal{F} \qquad \text{for every Borel set } B \subseteq \mathbb{R}.\]
    Its distribution (or law) is the probability measure $P_X(B) = P(X \in B)$ on $(\mathbb{R}, \mathcal{B}(\mathbb{R}))$, and its cumulative distribution function (CDF) is $F(x) = P(X \leq x)$.
\end{definition}

The distribution $P_X$ is a probability measure, since taking preimages commutes with complements and countable unions; it is called the image (or pushforward) of $P$ under $X$. To check that $X$ is measurable, it suffices to check $\{X \leq x\} \in \mathcal{F}$ for all $x$: the Borel sets $B$ with $\{X \in B\} \in \mathcal{F}$ form a $\sigma$-algebra, and if it contains the intervals $(-\infty, x]$, it contains all Borel sets. One can show that continuous functions of random variables, and sums, products, maxima and limits of sequences of random variables, are again random variables. We use such facts without comment: in practice, every function that we write down is measurable.

\begin{theorem}[The CDF determines the distribution]\label{thm:cdf}
    The CDF $F$ of a random variable $X$ is nondecreasing and right-continuous, with $\lim_{x \to -\infty} F(x) = 0$ and $\lim_{x \to \infty} F(x) = 1$, and $P(a < X \leq b) = F(b) - F(a)$ for $a < b$. Random variables with the same CDF have the same distribution. Conversely, every function $F$ with these four properties is the CDF of a random variable.
\end{theorem}

\begin{proof}[Proof of some parts]
    $F$ is nondecreasing by part (2) of Proposition~\ref{L4-prop:probability-rules} of Lecture~4. If $x_k \downarrow x$, the events $\{X \leq x_k\}$ decrease to $\{X \leq x\}$, so $F(x_k) \to F(x)$ by continuity (Proposition~\ref{prop:continuity}); the limits at $\pm\infty$ follow in the same way from $\bigcap_k \{X \leq -k\} = \emptyset$ and $\bigcup_k \{X \leq k\} = \Omega$. That the CDF determines $P_X$ follows from a uniqueness theorem of measure theory, which we do not prove: two probability measures that agree on a collection of sets that is closed under intersections, such as the intervals $(-\infty, x]$, agree on the $\sigma$-algebra that it generates. The converse follows from inverse transform sampling (Theorem~\ref{L4-thm:inverse-transform} of Lecture~4): on the probability space $\big((0, 1), \mathcal{B}((0, 1)), \lambda\big)$, the random variable $U(\omega) = \omega$ is uniform, and $F^{-1}(U)$ has the CDF $F$.
\end{proof}

So in practice, we rarely look at the sample space: a random variable is described by its distribution, i.e., by its CDF. The jumps of the CDF are the values that have positive probability: since $\{X < x\} = \bigcup_k \{X \leq x - \frac{1}{k}\}$, continuity gives $P(X < x) = \lim_{y \uparrow x} F(y)$, so $P(X = x) = F(x) - \lim_{y \uparrow x} F(y)$.

The discrete and continuous random variables of Definition~\ref{L4-def:pmf-pdf} of Lecture~4 are the two most common kinds. The density $p$ of a continuous random variable has to be measurable, and then $P(X \in B) = \int_B p(x)\, dx$ for every Borel set $B$, since both sides are probability measures that agree on intervals. Here, and from now on, integrals are Lebesgue integrals (Section~\ref{sec:lebesgue}), which agree with the integrals of Lecture~3 whenever these converge absolutely. In Section~\ref{sec:densities}, we will see that pmfs and pdfs are both densities with respect to a measure, and that some distributions have neither.

\section{The Lebesgue Integral and Expectation}\label{sec:lebesgue}

The expectation of a random variable is its average value, weighted by probability. For a random variable with finitely many values $x_1, \ldots, x_k$, it is the sum $\sum_i x_iP(X = x_i)$. In general, it is an integral with respect to the probability measure $P$, the Lebesgue integral. The construction works for every measure, and for the Lebesgue measure on $\mathbb{R}^n$, it extends the integral of Lecture~3.

Let $(\Omega, \mathcal{F}, \mu)$ be a measure space. A function $f: \Omega \to \mathbb{R}$ or $f: \Omega \to [0, \infty]$ is measurable if $\{f \leq x\} \in \mathcal{F}$ for all $x \in \mathbb{R}$; on a probability space, the measurable real-valued functions are the random variables.

\begin{definition}[Lebesgue integral]\label{def:lebesgue}
    \begin{enumerate}
        \item A simple function is a finite sum $s = \sum_{i=1}^k a_i\mathbf{1}_{A_i}$ with $a_i \geq 0$ and $A_i \in \mathcal{F}$, and its integral is $\int s\, d\mu = \sum_{i=1}^k a_i\mu(A_i)$, with $0 \cdot \infty = 0$.
        \item For a measurable $f: \Omega \to [0, \infty]$,
        \[\int f\, d\mu = \sup\Big\{\int s\, d\mu : s \text{ simple, } 0 \leq s \leq f\Big\} \in [0, \infty].\]
        \item A measurable $f: \Omega \to \mathbb{R}$ is integrable if $\int |f|\, d\mu < \infty$, and then $\int f\, d\mu = \int f^+\, d\mu - \int f^-\, d\mu$, where $f^+ = \max(f, 0)$ and $f^- = \max(-f, 0)$, so that $f = f^+ - f^-$ and $|f| = f^+ + f^-$.
    \end{enumerate}
    For $A \in \mathcal{F}$, we write $\int_A f\, d\mu = \int f\mathbf{1}_A\, d\mu$. Other notations are $\int f(\omega)\, \mu(d\omega)$, and for the Lebesgue measure, $\int f(\mathbf{x})\, d\mathbf{x}$.
\end{definition}

(One checks that the integral of a simple function does not depend on how it is written as a sum.) The Riemann integral of Lecture~3 divides the domain into small boxes. The Lebesgue integral divides the range instead: the simple functions
\[s_n = \min\big(n, \, 2^{-n}\lfloor 2^nf\rfloor\big),\]
which round $f$ down to a multiple of $2^{-n}$ and cut it off at $n$, increase to $f$, and $\int s_n\, d\mu$ adds up the values $j2^{-n}$, weighted by the measures of the sets $\{j2^{-n} \leq f < (j + 1)2^{-n}\}$. One can show that the integral is linear and monotone ($f \leq g$ implies $\int f\, d\mu \leq \int g\, d\mu$), that $\int |f|\, d\mu = 0$ if and only if $f = 0$ almost everywhere, i.e., outside a set of measure $0$, and that functions which agree almost everywhere have the same integral. For bounded Riemann integrable functions on boxes, and for absolutely convergent improper Riemann integrals, the Lebesgue integral agrees with the Riemann integral. But it integrates many more functions. For example, the indicator function $\mathbf{1}_{\mathbb{Q}}$ of the rationals on $[0, 1]$ is not Riemann integrable, since every Riemann sum can be made $0$ or $1$ by choosing irrational or rational points, but its Lebesgue integral is $\lambda(\mathbb{Q} \cap [0, 1]) = 0$. For the counting measure on $\mathbb{N}$, the integral is a series: $\int f\, d\mu = \sum_k f(k)$.

\begin{definition}[Expectation]\label{def:expectation}
    The expectation (expected value, or mean) of a random variable $X$ on a probability space $(\Omega, \mathcal{F}, P)$ is
    \[\E[X] = \int X\, dP,\]
    defined if $X \geq 0$ (then possibly $\E[X] = \infty$) or if $X$ is integrable, i.e., $\E[|X|] < \infty$.
\end{definition}

For discrete random variables and random variables with a density, this is the expectation of Definition~\ref{L4-def:expectation} of Lecture~4, by Theorem~\ref{thm:lotus} below. The expectation is linear, $\E[aX + bY + c] = a\E[X] + b\E[Y] + c$, and monotone, since the integral is, and $\E[\mathbf{1}_A] = P(A)$ for every event $A$. So linearity holds for all integrable random variables; Lecture~4 only proved it for discrete ones. To compute expectations, we do not need the sample space, only the distribution.

\begin{theorem}[Expectations from distributions]\label{thm:lotus}
    Let $X$ be a random variable, and let $g: \mathbb{R} \to \mathbb{R}$ be measurable with $g \geq 0$ or $\E[|g(X)|] < \infty$. Then $\E[g(X)] = \int_{\mathbb{R}} g\, dP_X$. In particular,
    \[\E[g(X)] = \sum_i g(x_i)\,p(x_i) \qquad \text{or} \qquad \E[g(X)] = \int_{-\infty}^\infty g(x)\,p(x) \, dx\]
    for discrete $X$ with the pmf $p$ and continuous $X$ with the density $p$, respectively.
\end{theorem}

\begin{proof}
    For $g = \mathbf{1}_B$, both sides are $P(X \in B)$, by the definition of $P_X$. By linearity, the formula holds for simple functions $g$, and by the monotone convergence theorem (Theorem~\ref{thm:monotone}), applied to simple functions that increase to $g$, for all $g \geq 0$. For integrable $g(X)$, apply it to $g^+$ and $g^-$. The two formulas follow in the same way from Definition~\ref{def:density}, with the counting measure and the Lebesgue measure.
\end{proof}

\paragraph{$L^p$ spaces.} For $p \geq 1$, the random variables with $\E[|X|^p] < \infty$ form a vector space $L^p$, and $\norm{X}_p = \E[|X|^p]^{1/p}$ is a norm on it, if we do not distinguish random variables that are equal almost surely (the triangle inequality is Minkowski's inequality, which we do not prove). By Jensen's inequality (Theorem~\ref{L4-thm:jensen} of Lecture~4) for the convex function $y \mapsto y^{q/p}$ on $[0, \infty)$, $\norm{X}_p \leq \norm{X}_q$ for $p \leq q$, so $L^q \subseteq L^p$. The space $L^2$, with the inner product $\E[XY]$ of Lecture~4, is complete, which makes it a Hilbert space (see the appendix to Lecture~3); we use this for the conditional expectation in Section~\ref{sec:conditional-expectation}.

\section{Convergence Theorems}\label{sec:convergence-theorems}

When can we exchange a limit and an expectation? This question comes up constantly: for limits of random variables, for infinite series, and for derivatives of expectations with respect to parameters. The answer is the main advantage of the Lebesgue integral.

\begin{theorem}[Monotone convergence]\label{thm:monotone}
    Let $0 \leq f_1 \leq f_2 \leq \cdots$ be measurable functions with $f_n \to f$ pointwise. Then $\int f_n\, d\mu \to \int f\, d\mu$.
\end{theorem}

\begin{proof}
    By monotonicity, the integrals $\int f_n\, d\mu$ increase to a limit $L \leq \int f\, d\mu$. Conversely, let $s = \sum_i a_i\mathbf{1}_{A_i}$ be simple with $0 \leq s \leq f$, and let $0 < c < 1$. The sets $E_n = \{f_n \geq cs\}$ increase to $\Omega$: where $s > 0$, we have $f_n \to f \geq s > cs$, and where $s = 0$, trivially $f_n \geq cs$. Hence
    \[\int f_n\, d\mu \geq \int cs\mathbf{1}_{E_n}\, d\mu = c\sum_i a_i\mu(A_i \cap E_n) \to c\int s\, d\mu,\]
    by the continuity of measures (Proposition~\ref{prop:continuity}). So $L \geq c\int s\, d\mu$ for all $c < 1$ and all such $s$, i.e., $L \geq \int f\, d\mu$.
\end{proof}

Applied to partial sums, monotone convergence shows that the expectation of a series of nonnegative random variables is the series of their expectations: $\E\big[\sum_k X_k\big] = \sum_k \E[X_k]$ for $X_k \geq 0$. For example, a random variable $X$ with values in $\{0, 1, 2, \ldots\}$ is the series $\sum_{k \geq 0} \mathbf{1}_{\{X > k\}}$, which gives the tail sum formula $\E[X] = \sum_{k \geq 0} P(X > k)$ of Example~\ref{L4-ex:tail-sum} of Lecture~4.

\begin{theorem}[Fatou's lemma and dominated convergence]\label{thm:dominated}
    \begin{enumerate}
        \item (Fatou) Let $f_n \geq 0$ be measurable. Then $\int \liminf_n f_n\, d\mu \leq \liminf_n \int f_n\, d\mu$.
        \item (Dominated convergence) Let $f_n \to f$ pointwise (or almost everywhere), and let $|f_n| \leq g$ for all $n$, with an integrable function $g$. Then $f$ is integrable, and $\int f_n\, d\mu \to \int f\, d\mu$.
    \end{enumerate}
\end{theorem}

\begin{proof}
    (1) The functions $g_k = \inf_{n \geq k} f_n$ increase to $\liminf_n f_n$, and $\int g_k\, d\mu \leq \int f_n\, d\mu$ for all $n \geq k$. By monotone convergence, $\int \liminf_n f_n\, d\mu = \lim_k \int g_k\, d\mu \leq \liminf_n \int f_n\, d\mu$. (2) Since $|f| \leq g$, $f$ is integrable. Fatou's lemma, applied to the nonnegative functions $g + f_n$ and $g - f_n$, gives
    \[\int g\, d\mu + \int f\, d\mu \leq \int g\, d\mu + \liminf_n \int f_n\, d\mu \qquad \text{and} \qquad \int g\, d\mu - \int f\, d\mu \leq \int g\, d\mu - \limsup_n \int f_n\, d\mu.\]
    Hence $\limsup_n \int f_n\, d\mu \leq \int f\, d\mu \leq \liminf_n \int f_n\, d\mu$.
\end{proof}

\begin{example}\label{ex:escaping-mass}
    \begin{enumerate}
        \item (Mass can escape.) On $[0, 1]$ with the Lebesgue measure, let $f_n = n\mathbf{1}_{(0, 1/n)}$. Then $f_n(x) \to 0$ for every $x$, but $\int f_n\, d\lambda = 1$ for all $n$: the limit of the integrals is not the integral of the limit. In the language of probability, $X_n = n\mathbf{1}_{\{U < 1/n\}}$ for a uniform $U$ is a lottery ticket that wins $n$ euros with probability $\frac{1}{n}$; it is $0$ for all large $n$, but $\E[X_n] = 1$. Dominated convergence does not apply, since $\sup_n f_n(x) \geq \frac{1}{2x}$ for $0 < x \leq \frac{1}{2}$, which is not integrable.
        \item (Differentiating under the integral sign.) Let $f(x, \theta)$ be integrable in $x$ and differentiable in $\theta$, with $|\frac{\partial f}{\partial\theta}(x, \theta)| \leq g(x)$ for all $\theta$ in an open interval, where $\int g\, d\mu < \infty$. Then
        \[\frac{d}{d\theta}\int f(x, \theta)\, d\mu(x) = \int \frac{\partial f}{\partial\theta}(x, \theta)\, d\mu(x).\]
        Indeed, for $h_n \to 0$, the difference quotients $\frac{f(x, \theta + h_n) - f(x, \theta)}{h_n}$ converge to $\frac{\partial f}{\partial\theta}(x, \theta)$, and by the mean value theorem, they are bounded by $g(x)$; now apply dominated convergence. This proves the Leibniz rule (Proposition~\ref{L3-prop:leibniz} of Lecture~3) for unbounded regions, under the condition stated there, and for general measures. It also shows that the gradient of an expected loss is the expected gradient, $\nabla_{\boldsymbol{\theta}}\E[\ell(\boldsymbol{\theta}; \mathbf{X})] = \E[\nabla_{\boldsymbol{\theta}}\ell(\boldsymbol{\theta}; \mathbf{X})]$, if the partial derivatives are dominated in this way. This is why the gradient on a random training example is an unbiased estimate of the gradient of the expected loss.
    \end{enumerate}
\end{example}

\section{Densities and the Radon--Nikodym Theorem}\label{sec:densities}

Discrete and continuous distributions look different, but they are two cases of one notion.

\begin{definition}[Density with respect to a measure]\label{def:density}
    Let $\mu$ be a measure on $(\mathbb{R}^n, \mathcal{B}(\mathbb{R}^n))$. A random vector $\mathbf{X}$, or its distribution, has the density $p \geq 0$ with respect to $\mu$ if
    \[P(\mathbf{X} \in B) = \int_B p\, d\mu \qquad \text{for all Borel sets } B.\]
    Then $\E[g(\mathbf{X})] = \int g\,p\, d\mu$ for every measurable $g \geq 0$, by the argument of Theorem~\ref{thm:lotus}.
\end{definition}

For the Lebesgue measure, $p$ is the density of Definition~\ref{L4-def:pmf-pdf} of Lecture~4. For the counting measure on a countable set $\{x_1, x_2, \ldots\}$ of values, $\int_B p\, d\mu = \sum_{x_i \in B} p(x_i)$, so $p$ is the pmf. So sums with pmfs and integrals with pdfs are integrals with respect to different reference measures, and everything that we prove for integrals holds for both. A necessary condition for a density is that $\mu(B) = 0$ implies $P(\mathbf{X} \in B) = 0$, since an integral over a null set is $0$. The distribution is then called absolutely continuous with respect to $\mu$, and this condition is also sufficient.

\begin{theorem}[Radon--Nikodym]\label{thm:radon-nikodym}
    Let $\mu$ and $\nu$ be $\sigma$-finite measures on $(\Omega, \mathcal{F})$ such that $\mu(A) = 0$ implies $\nu(A) = 0$ (written $\nu \ll \mu$). Then there is a measurable $p \geq 0$ with $\nu(A) = \int_A p\, d\mu$ for all $A \in \mathcal{F}$, and $p$ is unique up to changes on a set of $\mu$-measure $0$. It is called the Radon--Nikodym derivative $\frac{d\nu}{d\mu}$.
\end{theorem}

We do not prove this theorem here; the appendix to Lecture~3 proves it with the Riesz representation theorem for Hilbert spaces. Not every distribution has a density with respect to the Lebesgue measure or a counting measure.

\begin{example}[The ReLU of a Gaussian]\label{ex:relu-gaussian}
    Let $Z \sim \mathcal{N}(0, 1)$, and let $Y = \max(0, Z)$ be the activation of a ReLU unit with a Gaussian input. Then $P(Y = 0) = P(Z \leq 0) = \frac{1}{2}$, and $P(a \leq Y \leq b) = \int_a^b \varphi(y)\, dy$ for $0 < a < b$. So $Y$ is neither discrete nor continuous: it has an atom of mass $\frac{1}{2}$ at $0$, which the Lebesgue measure does not see. But $Y$ has a density with respect to $\mu = \delta_0 + \lambda$, namely $p(0) = \frac{1}{2}$ and $p(y) = \varphi(y)\mathbf{1}_{\{y > 0\}}$ for $y \neq 0$. Hence
    \[\E[Y] = 0 \cdot \tfrac{1}{2} + \int_0^\infty y\varphi(y)\, dy = \frac{1}{\sqrt{2\pi}} \approx 0.399, \qquad \E[Y^2] = \int_0^\infty y^2\varphi(y)\, dy = \frac{1}{2},\]
    so $\Var(Y) = \frac{1}{2} - \frac{1}{2\pi} \approx 0.341$. Mixed distributions of this kind also describe zero-inflated data, such as daily rainfall, or the amount that a visitor of a web shop spends.
\end{example}

\begin{remark}[Singular distributions and likelihood ratios]\label{rem:singular}
    \begin{enumerate}
        \item Some distributions have neither atoms nor a density with respect to the Lebesgue measure. An example is the Cantor distribution, the distribution of $X = \sum_{k=1}^\infty 2D_k3^{-k}$ with independent fair coin flips $D_k \in \{0, 1\}$, a random number whose ternary digits are all $0$ or $2$. Its CDF, the Cantor function, is continuous, so $X$ has no atoms. But $X$ lies in the Cantor set, which has Lebesgue measure $0$, so it has no density. One can show that every distribution on $\mathbb{R}$ is a mixture of a discrete, an absolutely continuous, and such a singular part (the Lebesgue decomposition).
        \item If $P \ll Q$, the Radon--Nikodym derivative $\frac{dP}{dQ}$ is the likelihood ratio, which is $\frac{p}{q}$ if $P$ and $Q$ have the densities $p$ and $q$. The weights of importance sampling (Section~\ref{L4-sec:monte-carlo} of Lecture~4) are likelihood ratios, $\E_P[f] = \E_Q[f\frac{dP}{dQ}]$, and the KL divergence of Section~\ref{L4-sec:information} of Lecture~4 is $\KL(P \,\|\, Q) = \E_P[\log\frac{dP}{dQ}]$ for $P \ll Q$, and $\infty$ otherwise. So these notions make sense for arbitrary distributions.
    \end{enumerate}
\end{remark}

\section{Product Measures and Independence}\label{sec:product}

A random vector $\mathbf{X} = (X_1, \ldots, X_n)$ is a measurable map from $\Omega$ to $\mathbb{R}^n$, i.e., $\{\mathbf{X} \in B\} \in \mathcal{F}$ for all Borel sets $B \subseteq \mathbb{R}^n$; one can show that this holds if and only if every $X_i$ is a random variable. Its joint distribution is the probability measure $P_{\mathbf{X}}(B) = P(\mathbf{X} \in B)$ on $\mathbb{R}^n$, and the distributions of the $X_i$ are its marginals. Joint pmfs and joint densities (Definition~\ref{L4-def:joint} of Lecture~4) are densities of the joint distribution with respect to the counting measure and the Lebesgue measure, and the ``reasonable'' sets in that definition are the Borel sets. On a product of two measure spaces, there is a natural measure.

\begin{theorem}[Product measures and Fubini--Tonelli]\label{thm:fubini-tonelli}
    Let $\mu$ and $\nu$ be $\sigma$-finite measures on $(\Omega_1, \mathcal{F}_1)$ and $(\Omega_2, \mathcal{F}_2)$, and let $\mathcal{F}_1 \otimes \mathcal{F}_2$ be the $\sigma$-algebra on $\Omega_1 \times \Omega_2$ generated by the products $A \times B$ with $A \in \mathcal{F}_1$ and $B \in \mathcal{F}_2$. There is exactly one measure $\mu \otimes \nu$ on it with $(\mu \otimes \nu)(A \times B) = \mu(A)\nu(B)$, the product measure. For every measurable $f \geq 0$ (Tonelli), and for every $f$ that is integrable with respect to $\mu \otimes \nu$ (Fubini),
    \[\int f\, d(\mu \otimes \nu) = \int\bigg(\int f(x, y)\, d\nu(y)\bigg) d\mu(x) = \int\bigg(\int f(x, y)\, d\mu(x)\bigg) d\nu(y).\]
\end{theorem}

We do not prove this theorem. The Lebesgue measure on $\mathbb{R}^2$ is $\lambda \otimes \lambda$, and the theorem contains Fubini's theorem of Lecture~3, including its condition: for nonnegative functions, the order of integration never matters, and for functions of both signs, it does not matter if $\int |f| < \infty$ (compare with Example~\ref{L3-ex:fubini} of Lecture~3, in which the two orders give $\frac{\pi}{4}$ and $-\frac{\pi}{4}$).

\begin{definition}[Independent random variables]\label{def:independent-rv}
    Random variables $X_1, \ldots, X_n$ are independent if
    \[P(X_1 \in B_1, \ldots, X_n \in B_n) = P(X_1 \in B_1)\cdots P(X_n \in B_n) \qquad \text{for all Borel sets } B_1, \ldots, B_n,\]
    i.e., if their joint distribution is the product $P_{X_1} \otimes \cdots \otimes P_{X_n}$ of their distributions. Infinitely many random variables are independent if any finitely many of them are.
\end{definition}

For discrete random variables, independence means $p(x, y) = p_X(x)\,p_Y(y)$ for all values $x$ and $y$, and for random variables with a joint density, it means that the joint density factorizes as $p(x, y) = p_X(x)\,p_Y(y)$ (outside a null set), as in Definition~\ref{L4-def:independent-rv} of Lecture~4. In the language of $\sigma$-algebras, the information contained in $X$ is the $\sigma$-algebra $\sigma(X) = \{\{X \in B\} : B \text{ Borel}\}$ generated by $X$, and $X$ and $Y$ are independent if and only if $\sigma(X)$ and $\sigma(Y)$ are independent (Definition~\ref{def:independence}). Since $\sigma(g(X)) \subseteq \sigma(X)$ for measurable $g$, functions $g(X)$ and $h(Y)$ of independent random variables are again independent.

\begin{proposition}\label{prop:independent-product}
    If $X$ and $Y$ are independent, and $g(X)$ and $h(Y)$ are integrable (or nonnegative), then $\E[g(X)h(Y)] = \E[g(X)]\,\E[h(Y)]$. In particular, $\E[XY] = \E[X]\,\E[Y]$ for integrable $X$ and $Y$.
\end{proposition}

\begin{proof}
    The joint distribution is $P_X \otimes P_Y$, so by Theorem~\ref{thm:lotus} and Tonelli's theorem (for nonnegative functions) or Fubini's theorem,
    \[\E[g(X)h(Y)] = \int\!\!\int g(x)h(y)\, dP_X(x)\, dP_Y(y) = \int g\, dP_X\int h\, dP_Y = \E[g(X)]\,\E[h(Y)]. \qedhere\]
\end{proof}

This extends Proposition~\ref{L4-prop:independent-product} of Lecture~4 to all random variables. Tonelli's theorem also proves the formulas of Examples~\ref{L4-ex:convolution} and~\ref{L4-ex:tail-integral} of Lecture~4 in general. For example, every random variable $X \geq 0$ satisfies $X = \int_0^\infty \mathbf{1}_{\{t < X\}}\, dt$, and Tonelli's theorem for $P \otimes \lambda$ gives
\[\E[X] = \int_0^\infty P(X > t)\, dt,\]
also if $X$ is neither discrete nor continuous.

\section{Conditional Expectation}\label{sec:conditional-expectation}

The conditional expectation $\E[Y \mid X]$ of Definition~\ref{L4-def:conditional-expectation} of Lecture~4 only depends on the information contained in $X$: it is a function of $X$, and by the proof of Theorem~\ref{L4-thm:conditional-expectation} of Lecture~4, with $h = \mathbf{1}_B$, $\E\big[(Y - \E[Y \mid X])\mathbf{1}_G\big] = 0$ for every event $G = \{X \in B\}$ in $\sigma(X)$. These two properties define the conditional expectation in general: also when $P(X = x) = 0$ for all $x$ and there is no joint density, and for conditioning on infinitely many random variables.

\begin{definition}[Conditional expectation given a $\sigma$-algebra]\label{def:conditional-expectation-general}
    Let $Y$ be integrable, and let $\mathcal{G} \subseteq \mathcal{F}$ be a $\sigma$-algebra. A conditional expectation of $Y$ given $\mathcal{G}$ is a $\mathcal{G}$-measurable integrable random variable $\E[Y \mid \mathcal{G}]$ with
    \[\E\big[\E[Y \mid \mathcal{G}]\,\mathbf{1}_G\big] = \E[Y\mathbf{1}_G] \qquad \text{for all } G \in \mathcal{G}.\]
    For a random variable or random vector $X$, $\E[Y \mid X] = \E[Y \mid \sigma(X)]$.
\end{definition}

\begin{theorem}\label{thm:conditional-expectation-general}
    The conditional expectation exists, and it is unique up to changes on a null set. If $\E[Y^2] < \infty$, then $\E[Y \mid \mathcal{G}]$ is the orthogonal projection of $Y$ onto the subspace $L^2(\mathcal{G})$ of $\mathcal{G}$-measurable random variables in $L^2$: it minimizes $\E[(Y - Z)^2]$ over all $Z \in L^2(\mathcal{G})$. Moreover, for integrable $Y$:
    \begin{enumerate}
        \item (Tower property) $\E\big[\E[Y \mid \mathcal{G}]\big] = \E[Y]$, and $\E\big[\E[Y \mid \mathcal{G}] \mid \mathcal{H}\big] = \E[Y \mid \mathcal{H}]$ for $\sigma$-algebras $\mathcal{H} \subseteq \mathcal{G}$;
        \item (Taking out what is known) $\E[ZY \mid \mathcal{G}] = Z\,\E[Y \mid \mathcal{G}]$ for every bounded $\mathcal{G}$-measurable random variable $Z$;
        \item (Independence) if $\sigma(Y)$ and $\mathcal{G}$ are independent, then $\E[Y \mid \mathcal{G}] = \E[Y]$.
    \end{enumerate}
\end{theorem}

\begin{proof}[Proof of some parts]
    Existence for $Y \in L^2$: one can show that $L^2(\mathcal{G})$ is a closed subspace of $L^2$, so by the projection theorem for Hilbert spaces (see the appendix to Lecture~3), $Y$ has an orthogonal projection $\hat Y$ onto it. The error $Y - \hat Y$ is orthogonal to every element of $L^2(\mathcal{G})$, in particular to $\mathbf{1}_G$ for $G \in \mathcal{G}$, which is the defining property. For integrable $Y \geq 0$, one applies the Radon--Nikodym theorem to the measure $G \mapsto \E[Y\mathbf{1}_G]$ on $\mathcal{G}$. Uniqueness: if $Z_1$ and $Z_2$ both have the defining property, then $G = \{Z_1 > Z_2\} \in \mathcal{G}$ gives $\E[(Z_1 - Z_2)\mathbf{1}_G] = 0$, so $Z_1 \leq Z_2$ almost surely, and by symmetry, $Z_1 = Z_2$ almost surely. (1) The first statement is the defining property for $G = \Omega$. For the second, $\E[Y \mid \mathcal{H}]$ is $\mathcal{H}$-measurable, and for $H \in \mathcal{H} \subseteq \mathcal{G}$, $\E\big[\E[Y \mid \mathcal{H}]\mathbf{1}_H\big] = \E[Y\mathbf{1}_H] = \E\big[\E[Y \mid \mathcal{G}]\mathbf{1}_H\big]$, which is the defining property of $\E\big[\E[Y \mid \mathcal{G}] \mid \mathcal{H}\big]$. (3) The constant $\E[Y]$ is $\mathcal{G}$-measurable, and $\E[Y\mathbf{1}_G] = \E[Y]\,P(G) = \E\big[\E[Y]\mathbf{1}_G\big]$ by independence (Proposition~\ref{prop:independent-product}). We omit the proof of (2).
\end{proof}

By the Doob--Dynkin lemma, which we do not prove, every $\sigma(X)$-measurable random variable is a measurable function of $X$. So $\E[Y \mid X] = m(X)$ for a measurable function $m$, and for discrete $X$ or a joint density, $m(x)$ is the conditional mean $\E[Y \mid X = x]$ of Definition~\ref{L4-def:conditional-expectation} of Lecture~4; Theorem~\ref{L4-thm:conditional-expectation} of Lecture~4 is the case $\mathcal{G} = \sigma(X)$ of Theorem~\ref{thm:conditional-expectation-general}.

\begin{example}\label{ex:conditional-partition}
    \begin{enumerate}
        \item (Finitely many pieces of information.) Let $\mathcal{G}$ be generated by a partition $B_1, \ldots, B_n$ of $\Omega$ with $P(B_i) > 0$ (Example~\ref{ex:sigma-algebras}). A $\mathcal{G}$-measurable random variable is constant on each $B_i$, and the defining property with $G = B_i$ gives
        \[\E[Y \mid \mathcal{G}] = \sum_{i=1}^n \frac{\E[Y\mathbf{1}_{B_i}]}{P(B_i)}\,\mathbf{1}_{B_i}:\]
        on each piece, $\E[Y \mid \mathcal{G}]$ is the average of $Y$ over that piece.
        \item (Law of total variance.) For $Y \in L^2$, the decomposition $Y - \E[Y] = \big(Y - \E[Y \mid \mathcal{G}]\big) + \big(\E[Y \mid \mathcal{G}] - \E[Y]\big)$ is orthogonal, and Pythagoras gives
        \[\Var(Y) = \E\big[\Var(Y \mid \mathcal{G})\big] + \Var\big(\E[Y \mid \mathcal{G}]\big), \qquad \text{where } \Var(Y \mid \mathcal{G}) = \E\big[(Y - \E[Y \mid \mathcal{G}])^2 \mid \mathcal{G}\big].\]
        For $\mathcal{G} = \sigma(X)$, this is the law of total variance of Example~\ref{L4-ex:total-variance} of Lecture~4.
    \end{enumerate}
\end{example}

\section{Martingales}\label{sec:martingales}

Conditional expectations describe what we expect, given the information that is available at a certain time. A filtration is an increasing sequence of $\sigma$-algebras $\mathcal{F}_0 \subseteq \mathcal{F}_1 \subseteq \cdots \subseteq \mathcal{F}$, where $\mathcal{F}_n$ is the information available at time $n$; for example, $\mathcal{F}_n = \sigma(X_1, \ldots, X_n)$ for a sequence of observations $X_1, X_2, \ldots$, and $\mathcal{F}_0 = \{\emptyset, \Omega\}$.

\begin{definition}[Martingale]\label{def:martingale}
    A sequence of integrable random variables $M_0, M_1, \ldots$ is a martingale with respect to a filtration $(\mathcal{F}_n)$ if every $M_n$ is $\mathcal{F}_n$-measurable and $\E[M_{n+1} \mid \mathcal{F}_n] = M_n$ for all $n$.
\end{definition}

A martingale is a fair game: given the past, the expected next value is the current value. By the tower property, $\E[M_n] = \E[M_0]$ for all $n$.

\begin{example}\label{ex:martingales}
    \begin{enumerate}
        \item (Random walks.) Let $X_1, X_2, \ldots$ be independent with mean $0$, and let $S_n = X_1 + \cdots + X_n$, with $S_0 = 0$. Taking out what is known and independence give $\E[S_{n+1} \mid \mathcal{F}_n] = S_n + \E[X_{n+1}] = S_n$, so $(S_n)$ is a martingale. If the $X_i$ also have the variance $\sigma^2$, then $S_n^2 - n\sigma^2$ is a martingale, since $\E[S_{n+1}^2 \mid \mathcal{F}_n] = S_n^2 + 2S_n\E[X_{n+1}] + \E[X_{n+1}^2] = S_n^2 + \sigma^2$.
        \item (Accumulating information.) For every integrable $Y$, $M_n = \E[Y \mid \mathcal{F}_n]$ is a martingale, by the tower property: it is the best prediction of $Y$, updated as information arrives. For example, the posterior mean $\E[\theta \mid X_1, \ldots, X_n]$ of Bayesian inference (Section~\ref{L4-sec:bayesian} of Lecture~4) is a martingale.
        \item (Likelihood ratios.) Let $X_1, X_2, \ldots$ be i.i.d.\ with the density $q$, and let $p$ be another density with $p(x) = 0$ whenever $q(x) = 0$. Then the likelihood ratio $L_n = \prod_{i=1}^n \frac{p(X_i)}{q(X_i)}$ of the data under $p$ versus $q$ is a martingale with $\E[L_n] = 1$, since its factors are independent with the mean $\int \frac{p}{q}\,q\, dx = 1$. By Ville's inequality, which we do not prove, $P\big(L_n \geq \frac{1}{\alpha} \text{ for some } n\big) \leq \alpha$ for every nonnegative martingale with $L_0 = 1$. So we may compute $L_n$ after every observation and reject the hypothesis that the data have the density $q$ as soon as $L_n \geq \frac{1}{\alpha}$: unlike repeated z-tests (Section~\ref{L4-sec:evaluation} of Lecture~4), this sequential test keeps its level $\alpha$, however often we look at the data. If the data actually have the density $p$, the strong law of large numbers (Theorem~\ref{thm:slln}) gives $\frac{1}{n}\log L_n \to \E_p\big[\log\frac{p(X_1)}{q(X_1)}\big] = \KL(p \,\|\, q)$ (Section~\ref{L4-sec:information} of Lecture~4), so $L_n$ grows exponentially fast, and the test eventually rejects.
    \end{enumerate}
\end{example}

A stopping time is a random time $\tau$ with values in $\{0, 1, 2, \ldots\} \cup \{\infty\}$ such that $\{\tau = n\} \in \mathcal{F}_n$ for all $n$: whether we stop at time $n$ only depends on the information available at time $n$, not on the future.

\begin{theorem}[Optional stopping]\label{thm:optional-stopping}
    Let $(M_n)$ be a martingale and $\tau$ a stopping time with $\tau < \infty$ almost surely, and let $|M_{\min(n, \tau)}| \leq C$ for all $n$, with a constant $C$. Then $\E[M_\tau] = \E[M_0]$.
\end{theorem}

\begin{proof}
    The stopped process $N_n = M_{\min(n, \tau)}$ is again a martingale: $N_{n+1} - N_n = \mathbf{1}_{\{\tau > n\}}(M_{n+1} - M_n)$, and since $\{\tau > n\} \in \mathcal{F}_n$, taking out what is known gives $\E[N_{n+1} - N_n \mid \mathcal{F}_n] = \mathbf{1}_{\{\tau > n\}}\E[M_{n+1} - M_n \mid \mathcal{F}_n] = 0$. So $\E[N_n] = \E[M_0]$ for all $n$. As $n \to \infty$, $N_n \to M_\tau$ almost surely, and $|N_n| \leq C$, so dominated convergence gives $\E[M_\tau] = \E[M_0]$.
\end{proof}

So no stopping rule turns a fair game into a favorable one, as long as the stakes stay bounded. Without this condition, the theorem fails: a gambler who doubles the stake after every loss in a fair game wins $1$ euro almost surely, but needs an unbounded fortune to do so.

\begin{example}[Gambler's ruin]\label{ex:gamblers-ruin}
    A gambler starts with $a$ euros and bets $1$ euro on fair coin flips, until the fortune reaches $a + b$ euros or $0$. The fortune is $a + S_n$ with the simple random walk $S_n$ (with steps $\pm 1$ of probability $\frac{1}{2}$ each), and the game ends at the first time $\tau$ with $S_\tau = b$ or $S_\tau = -a$. One can show that $\tau < \infty$ almost surely, and $|S_{\min(n, \tau)}| \leq \max(a, b)$. By optional stopping,
    \[0 = \E[S_\tau] = b\,P(S_\tau = b) - a\,P(S_\tau = -a), \qquad \text{so} \qquad P(S_\tau = b) = \frac{a}{a + b}.\]
    Applying optional stopping to the martingale $S_n^2 - n$ (with more care, since it is not bounded) gives $\E[\tau] = \E[S_\tau^2] = b^2\frac{a}{a + b} + a^2\frac{b}{a + b} = ab$. For $a = 10$ and $b = 90$, the gambler reaches $100$ euros with probability $0.1$, and the game lasts $900$ rounds on average.
\end{example}

Martingales also exist in continuous time. With the filtration $\mathcal{F}_s = \sigma(B(u) : u \leq s)$, the independence of the increments of Brownian motion (Definition~\ref{L4-def:brownian} of Lecture~4) gives $\E[B(t) \mid \mathcal{F}_s] = B(s) + \E[B(t) - B(s)] = B(s)$ for $s \leq t$: Brownian motion is a martingale, the continuous analogue of a fair random walk.

\section{Modes of Convergence and the Strong Law}\label{sec:convergence}

The weak law of large numbers (Theorem~\ref{L4-thm:lln} of Lecture~4) is a statement about probabilities, for each $N$. Measure theory allows statements about the whole random sequence: for example, that $\bar X_N(\omega) \to \mu$ for almost every outcome $\omega$, i.e., along almost every infinite sequence of observations. Random variables can converge in several different senses.

\begin{definition}[Modes of convergence]\label{def:convergence-modes}
    Let $X, X_1, X_2, \ldots$ be random variables on the same probability space. We say that $X_n \to X$
    \begin{enumerate}
        \item almost surely if $P(X_n \to X) = 1$, i.e., $X_n(\omega) \to X(\omega)$ for all $\omega$ outside a null set;
        \item in probability if $P(|X_n - X| > \varepsilon) \to 0$ for every $\varepsilon > 0$;
        \item in $L^p$, for $p \geq 1$, if $\E[|X_n - X|^p] \to 0$; for $p = 2$, this is convergence in mean square.
    \end{enumerate}
    Random variables $X_n$ and $X$, possibly on different probability spaces, converge in distribution if $F_{X_n}(x) \to F_X(x)$ at every point $x$ at which $F_X$ is continuous.
\end{definition}

Convergence in probability and in distribution are the notions of Definition~\ref{L4-def:convergence} of Lecture~4: the weak law of large numbers is convergence in probability, and the central limit theorem (Theorem~\ref{L4-thm:clt} of Lecture~4) is convergence in distribution of the standardized means to $\mathcal{N}(0, 1)$.

\begin{proposition}\label{prop:convergence-modes}
    Almost sure convergence and convergence in $L^p$ each imply convergence in probability, which implies convergence in distribution. None of the converse implications holds in general.
\end{proposition}

\begin{proof}
    If $X_n \to X$ almost surely, the events $A_n = \{|X_m - X| > \varepsilon \text{ for some } m \geq n\}$ decrease to an event that is contained in the null event $\{X_n \not\to X\}$, so $P(|X_n - X| > \varepsilon) \leq P(A_n) \to 0$ by continuity. Convergence in $L^p$ implies convergence in probability by Markov's inequality (Proposition~\ref{L4-prop:markov} of Lecture~4): $P(|X_n - X| > \varepsilon) \leq \frac{\E[|X_n - X|^p]}{\varepsilon^p}$. If $X_n \to X$ in probability, then $X_n = X + (X_n - X)$ with $X_n - X \to 0$ in probability, so $X_n \to X$ in distribution by Slutsky's theorem (Theorem~\ref{L4-thm:slutsky} of Lecture~4), applied to the constant sequence $X$. Example~\ref{ex:convergence-counterexamples} disproves the converses.
\end{proof}

\begin{example}\label{ex:convergence-counterexamples}
    Let $U$ be uniform on $[0, 1)$.
    \begin{enumerate}
        \item The lottery tickets $X_n = n\mathbf{1}_{\{U < 1/n\}}$ of Example~\ref{ex:escaping-mass} converge to $0$ almost surely, since $X_n = 0$ as soon as $n > \frac{1}{U}$, but not in $L^1$, since $\E[X_n] = 1$.
        \item (A sliding block.) Write $n = 2^k + j$ with $0 \leq j < 2^k$, and let $X_n = \mathbf{1}_{\{j2^{-k} \leq U < (j + 1)2^{-k}\}}$: the indicator of a block of length $2^{-k}$, which slides across $[0, 1)$, and then starts again with half the length. Then $\E[|X_n|^p] = 2^{-k} \to 0$, so $X_n \to 0$ in every $L^p$ and in probability. But for every value of $U$, one block of every generation $k$ covers it, so $X_n = 1$ for infinitely many $n$, and the sequence $X_n(\omega)$ converges for no outcome $\omega$.
        \item If $X \sim \mathcal{N}(0, 1)$ and $X_n = -X$ for all $n$, then $X_n \to X$ in distribution, since $-X$ has the same distribution as $X$, but $|X_n - X| = 2|X|$ does not converge to $0$ in probability.
    \end{enumerate}
\end{example}

The main tool for proving almost sure convergence is the following lemma.

\begin{lemma}[Borel--Cantelli]\label{lem:borel-cantelli}
    Let $A_1, A_2, \ldots$ be events, and let $\{A_n \text{ i.o.}\}$ be the event that infinitely many of them happen (``infinitely often'').
    \begin{enumerate}
        \item If $\sum_n P(A_n) < \infty$, then $P(A_n \text{ i.o.}) = 0$.
        \item If the $A_n$ are independent and $\sum_n P(A_n) = \infty$, then $P(A_n \text{ i.o.}) = 1$.
    \end{enumerate}
\end{lemma}

\begin{proof}
    (1) For every $k$, $\{A_n \text{ i.o.}\} \subseteq \bigcup_{n \geq k} A_n$, which has probability at most $\sum_{n \geq k} P(A_n)$ by the union bound, and this tends to $0$ as $k \to \infty$. (2) The complement of $\{A_n \text{ i.o.}\}$ is the union over $k$ of the events $\bigcap_{n \geq k} A_n^c$. By independence and $1 - x \leq e^{-x}$, for all $m \geq k$,
    \[P\Big(\bigcap_{n=k}^m A_n^c\Big) = \prod_{n=k}^m \big(1 - P(A_n)\big) \leq \exp\Big(-\sum_{n=k}^m P(A_n)\Big),\]
    which tends to $0$ as $m \to \infty$. So every $\bigcap_{n \geq k} A_n^c$ is a null event, and so is their countable union.
\end{proof}

\begin{example}\label{ex:borel-cantelli}
    Consider infinitely many fair coin flips (Example~\ref{ex:coin-flips}). The events $A_k$ that the flips $2k - 1$ and $2k$ are both heads are independent and have the probability $\frac{1}{4}$, so by part (2), almost surely, two heads in a row appear infinitely often; the same argument works for every finite pattern, for example for the text of this appendix, encoded in bits. On the other hand, let $B_k$ be the event that the $\lceil 2\log_2 k\rceil$ flips starting with flip $k$ are all heads. Then $P(B_k) \leq k^{-2}$, which is summable, so by part (1), almost surely, only finitely many $B_k$ happen: the longest run of heads among the first $n$ flips is eventually shorter than about $2\log_2 n$.
\end{example}

\begin{theorem}[Strong law of large numbers]\label{thm:slln}
    Let $X_1, X_2, \ldots$ be i.i.d.\ and integrable with mean $\mu$. Then $\bar X_N \to \mu$ almost surely.
\end{theorem}

\begin{proof}[Proof if $X_1$ has a finite fourth moment]
    Let $Y_i = X_i - \mu$ and $S_N = Y_1 + \cdots + Y_N$. When we expand $S_N^4 = \sum_{i, j, k, l} Y_iY_jY_kY_l$, every term in which some $Y_i$ appears exactly once has expectation $0$, by independence. What remains are the $N$ terms $Y_i^4$ and the $3N(N - 1)$ terms $Y_i^2Y_j^2$ with $i \neq j$, so $\E[S_N^4] = N\E[Y_1^4] + 3N(N - 1)\sigma^4 \leq CN^2$ for a constant $C$. By Markov's inequality,
    \[P(|\bar X_N - \mu| \geq \varepsilon) = P(S_N^4 \geq N^4\varepsilon^4) \leq \frac{C}{N^2\varepsilon^4},\]
    which is summable. By the Borel--Cantelli lemma, almost surely, $|\bar X_N - \mu| \geq \varepsilon$ only for finitely many $N$. Applying this to $\varepsilon = \frac{1}{m}$ for $m = 1, 2, \ldots$ (countably many null events) gives $\bar X_N \to \mu$ almost surely.
\end{proof}

\begin{example}[Normal numbers]\label{ex:normal-numbers}
    The binary digits of a uniformly random $U \in [0, 1)$ are i.i.d.\ fair coin flips (Example~\ref{ex:coin-flips}). By the strong law, the fraction of $1$s among the first $N$ digits converges to $\frac{1}{2}$ almost surely: almost every number in $[0, 1)$ has asymptotically equally many $0$s and $1$s in its binary expansion (Borel's normal number theorem, 1909). Yet for natural constants such as $\sqrt{2} - 1$, it is not known whether they have this property. This is typical of probabilistic arguments: they show that almost all objects have a property, without exhibiting one.
\end{example}

By the strong law, applied to the $X_i$ and the $X_i^2$, the sample standard deviation $s_N$ also converges to $\sigma$ almost surely. Similarly, every estimator that is a continuous function of sample means, such as the maximum likelihood estimator $\frac{1}{\bar X_N}$ of the rate of an exponential distribution, converges to the true parameter almost surely: it is strongly consistent, which is stronger than consistency (Definition~\ref{L4-def:consistent} of Lecture~4).

\section*{Key Takeaways}
\addcontentsline{toc}{section}{Key Takeaways}
\begin{itemize}[leftmargin=*]
    \item Events form a $\sigma$-algebra, which is closed under complements and countable unions, and a probability is a countably additive measure with $P(\Omega) = 1$. Not every subset of $[0, 1]$ can have a length (Vitali), but every Borel set has one (the Lebesgue measure). Measures are continuous along increasing and decreasing sequences of sets, and null events, such as a rational uniform point, can happen. $\sigma$-algebras describe information, and independent $\sigma$-algebras describe pieces of information that tell us nothing about each other.
    \item Random variables are measurable functions, and their distributions are the image measures $P_X$. The CDF determines the distribution, and every nondecreasing, right-continuous function with the limits $0$ and $1$ is a CDF.
    \item The Lebesgue integral divides the range instead of the domain; it extends the integral of Lecture~3 and contains series as integrals with respect to the counting measure. The expectation is $\E[X] = \int X\, dP$, computed from the distribution as $\E[g(X)] = \int g\, dP_X$. Monotone and dominated convergence allow us to exchange limits and expectations, and to differentiate under the integral sign; without a dominating function, mass can escape.
    \item Pmfs and pdfs are densities with respect to the counting and the Lebesgue measure. A distribution has a density with respect to $\mu$ if and only if it is absolutely continuous with respect to $\mu$ (Radon--Nikodym). Mixed distributions such as the ReLU of a Gaussian have densities with respect to other measures, singular distributions such as the Cantor distribution have none, and likelihood ratios are Radon--Nikodym derivatives.
    \item Independence means that the joint distribution is the product measure of the marginals. By Fubini--Tonelli, nonnegative or integrable functions can be integrated in either order, so expectations of products of independent random variables factorize, and $\E[X] = \int_0^\infty P(X > t)\, dt$ for every $X \geq 0$.
    \item The conditional expectation $\E[Y \mid \mathcal{G}]$ is the $\mathcal{G}$-measurable random variable with the same integrals as $Y$ over the events of $\mathcal{G}$; for $Y \in L^2$, it is the orthogonal projection of $Y$ onto $L^2(\mathcal{G})$, the best mean-squared-error prediction of $Y$ from the information $\mathcal{G}$. It satisfies the tower property and the law of total variance, and it takes out what is known.
    \item A martingale is a fair game, $\E[M_{n+1} \mid \mathcal{F}_n] = M_n$; random walks, posterior means, likelihood ratios and Brownian motion are martingales. Stopping at a stopping time does not change the expectation if the stopped martingale stays bounded (optional stopping), which solves the gambler's ruin, and Ville's inequality gives likelihood ratio tests that stay valid when we look at the data repeatedly.
    \item Almost sure convergence and convergence in $L^p$ imply convergence in probability, which implies convergence in distribution, but not conversely. The Borel--Cantelli lemma proves almost sure statements, and the strong law of large numbers says that $\bar X_N \to \mu$ almost surely; for example, almost every number is normal.
\end{itemize}

\section*{Exercises}
\addcontentsline{toc}{section}{Exercises}
\begin{enumerate}[leftmargin=*]
    \item ($\sigma$-algebras.)
    \begin{enumerate}[label=(\alph*)]
        \item Show that the intersection of two $\sigma$-algebras on $\Omega$ is a $\sigma$-algebra, but that their union need not be one. (Hint: $\Omega = \{1, 2, 3\}$.)
        \item Let $\mathcal{G}$ be the $\sigma$-algebra on $\Omega = \{1, \ldots, 6\}^2$ generated by the events $\{\text{the sum is } s\}$ for two dice. How many elements does $\mathcal{G}$ have? Does it contain the event ``both dice show the same number''?
    \end{enumerate}
    \item Let $X$ be a random variable with the CDF $F$. Use the continuity of probability (Proposition~\ref{prop:continuity}) to show that $P(X > t) \to 0$ as $t \to \infty$, and that $P(X = x) = F(x) - \lim_{y \uparrow x} F(y)$. Conclude that $F$ has at most countably many jumps. (Hint: there are at most $n$ jumps of size larger than $\frac{1}{n}$.)
    \item (Limits and expectations.)
    \begin{enumerate}[label=(\alph*)]
        \item Use monotone convergence to show that $\lim_{n \to \infty}\int_0^n \big(1 - \frac{x}{n}\big)^n\, dx = 1$. (Hint: $(1 - \frac{x}{n})^n\mathbf{1}_{[0, n]}(x)$ increases to $e^{-x}$.)
        \item Let $X \geq 0$ with $\E[X] < \infty$. Show that $\E[X\mathbf{1}_{\{X > t\}}] \to 0$ as $t \to \infty$, and conclude that $nP(X > n) \to 0$.
        \item Find random variables $X_n \geq 0$ with $X_n \to 0$ almost surely and $\E[X_n] \to \infty$.
    \end{enumerate}
    \item A waiting time $Z \sim$ Exponential$(1)$ is cut off after $2$ minutes: $Y = \min(Z, 2)$. Find the CDF of $Y$ and a density of $Y$ with respect to $\lambda + \delta_2$, and compute $\E[Y]$ in two ways: with this density, and with the tail integral $\E[Y] = \int_0^\infty P(Y > t)\, dt$ of Section~\ref{sec:product}.
    \item Let $X_1$ and $X_2$ be independent and Uniform$(0, 1)$, and $M = \max(X_1, X_2)$.
    \begin{enumerate}[label=(\alph*)]
        \item Show that $\E[X_1\mathbf{1}_{\{M \leq t\}}] = \E\big[\frac{3}{4}M\mathbf{1}_{\{M \leq t\}}\big]$ for $0 \leq t \leq 1$. (One can show that this implies $\E[X_1 \mid M] = \frac{3}{4}M$, since the events $\{M \leq t\}$ generate $\sigma(M)$ and are closed under intersections.)
        \item Explain the result: given $M = m$, $X_1 = m$ with probability $\frac{1}{2}$, and otherwise $X_1$ is uniform on $[0, m]$. Check the tower property.
    \end{enumerate}
    \item (P\'olya's urn.) An urn contains one red and one blue ball. Repeatedly, a ball is drawn at random and put back together with another ball of the same color. Let $R_n$ be the fraction of red balls after $n$ draws. Show that $(R_n)$ is a martingale, and find $\E[R_n]$.
    \item Let $X_n$ be independent with $P(X_n = 1) = p_n$ and $P(X_n = 0) = 1 - p_n$.
    \begin{enumerate}[label=(\alph*)]
        \item For $p_n = \frac{1}{n}$, show that $X_n \to 0$ in probability and in $L^1$, but not almost surely.
        \item For $p_n = \frac{1}{n^2}$, show that $X_n \to 0$ almost surely.
    \end{enumerate}
\end{enumerate}

\section*{Programming Exercises}
\addcontentsline{toc}{section}{Programming Exercises}

% Inline code cannot be hyphenated, so allow a little extra stretch in this section.
\setlength{\emergencystretch}{3em}

As in Lecture~4, implement the methods yourself with NumPy, and use library routines (such as \py{rng.normal} or \py{scipy.stats}) only to check your results. Start each file with
\begin{lstlisting}
import numpy as np
import matplotlib.pyplot as plt
rng = np.random.default_rng(0)
\end{lstlisting}

\begin{enumerate}[label=\textbf{P\arabic*.}, leftmargin=*]
    \item \textbf{Distributions without densities.}
    \begin{itemize}
        \item Sample $10^5$ values of $\max(0, Z)$ for $Z \sim \mathcal{N}(0, 1)$, and check $P(Y = 0) = \frac{1}{2}$, $\E[Y] = \frac{1}{\sqrt{2\pi}}$ and $\Var(Y) = \frac{1}{2} - \frac{1}{2\pi}$ (Example~\ref{ex:relu-gaussian}).
        \item Sample from the Cantor distribution (Remark~\ref{rem:singular}) with $30$ random ternary digits $0$ or $2$, plot the empirical CDF (the Cantor function), and check the mean $\frac{1}{2}$ and the variance $\frac{1}{8}$. What fraction of the samples lies in $(\frac{1}{3}, \frac{2}{3})$, and in $(\frac{1}{9}, \frac{2}{9})$?
    \end{itemize}

    \item \textbf{Modes of convergence.}
    \begin{itemize}
        \item For the sequences (1) and (2) of Example~\ref{ex:convergence-counterexamples}, plot the paths $n \mapsto X_n$ for $n \leq 1000$ for a few values of $U$, and estimate $\E[X_n]$ and $P(X_n \neq 0)$ from $10^4$ values of $U$. Which sequence converges almost surely, which in probability, and which in $L^1$?
        \item Simulate independent events $A_n$ with $P(A_n) = \frac{1}{n}$ and with $P(A_n) = \frac{1}{n^2}$, for $n \leq 10^6$. In $20$ runs, count how many of the events happen, and record when the last one happens. Compare with the Borel--Cantelli lemma; for $P(A_n) = \frac{1}{n}$, compare the count with $\sum_{n \leq N}\frac{1}{n} \approx \log N$.
        \item (The strong law.) Plot $\bar X_N$ against $N \leq 10^4$ for $20$ independent sequences of Exponential$(1)$ variables, together with the curves $1 \pm 2/\sqrt{N}$. How does the spread of the paths shrink, and does every path converge?
    \end{itemize}

    \item \textbf{Martingales.}
    \begin{itemize}
        \item Simulate the gambler's ruin of Example~\ref{ex:gamblers-ruin} with $a = 10$ and $b = 90$, $10\,000$ times, and check the winning probability $0.1$ and the mean duration $900$. Then let the coin show heads with probability $0.49$. How do the two numbers change?
        \item (A test that allows peeking.) For coin flips with $H_0$: $\theta = \frac{1}{2}$ and the alternative $\theta = 0.6$, compute the likelihood ratio $L_n = \prod_{i \leq n}\frac{0.6^{x_i}0.4^{1 - x_i}}{0.5}$ after every flip, and reject $H_0$ as soon as $L_n \geq 20$ (Example~\ref{ex:martingales}). In $10\,000$ simulations of a fair coin with up to $1000$ flips, how often is $H_0$ rejected? How many flips are needed on average to reject $H_0$ when $\theta = 0.6$? Compare with the peeking z-test in the programming exercises of Lecture~4.
    \end{itemize}
\end{enumerate}

\end{document}
